% Paso 2: Plantilla de Análisis Léxico - VERSIÓN CORREGIDA
% Proyecto: Análisis Exegético Profundo
% Ministerio "La Gloria es del Señor"

\documentclass[11pt,a4paper]{article}
\usepackage[utf8]{inputenc}
\usepackage[spanish]{babel}
\usepackage[margin=2.5cm]{geometry}
\usepackage{graphicx}
\usepackage{fancyhdr}
\usepackage{xcolor}
\usepackage{tcolorbox}
\usepackage{enumitem}
\usepackage{tabularx}
\usepackage{multirow}
\usepackage{amsfonts}
\usepackage{fontspec}
\usepackage{titlesec}
\usepackage{array}

% Configuración de colores
\definecolor{forestgreen}{RGB}{27,67,50}
\definecolor{sagegreen}{RGB}{45,90,39}
\definecolor{olivegreen}{RGB}{64,145,108}
\definecolor{mintgreen}{RGB}{82,183,136}
\definecolor{lightmint}{RGB}{116,198,157}
\definecolor{softcream}{RGB}{216,243,220}
\definecolor{papergray}{RGB}{118,128,150}

% Configuración de encabezados
\pagestyle{fancy}
\fancyhf{}
\fancyhead[L]{\textcolor{forestgreen}{\textbf{Análisis Exegético Profundo}}}
\fancyhead[R]{\textcolor{papergray}{Paso 2: Análisis Léxico}}
\fancyfoot[C]{\textcolor{papergray}{\thepage}}
\fancyfoot[R]{\textcolor{papergray}{\scriptsize Ministerio "La Gloria es del Señor"}}

% Configuración de títulos
\titleformat{\section}
{\color{forestgreen}\Large\bfseries}
{\thesection}{1em}{}

\titleformat{\subsection}
{\color{sagegreen}\large\bfseries}
{\thesubsection}{1em}{}

% Configuración de cajas
\tcbuselibrary{skins}
\newtcolorbox{infobox}[1]{
    colback=softcream,
    colframe=olivegreen,
    fonttitle=\bfseries,
    title=#1,
    left=10pt,
    right=10pt,
    top=10pt,
    bottom=10pt
}

\newtcolorbox{checklistbox}{
    colback=white,
    colframe=mintgreen,
    left=15pt,
    right=15pt,
    top=10pt,
    bottom=10pt,
    boxrule=2pt
}

\newtcolorbox{lexiconbox}{
    colback=lightmint!20,
    colframe=sagegreen,
    left=10pt,
    right=10pt,
    top=8pt,
    bottom=8pt,
    boxrule=1.5pt
}

% Definir tipos de columna personalizados
\newcolumntype{L}[1]{>{\raggedright\arraybackslash}p{#1}}
\newcolumntype{C}[1]{>{\centering\arraybackslash}p{#1}}

\begin{document}

% Título principal
\begin{center}
    {\Huge\color{forestgreen}\textbf{PASO 2}}\\[0.5cm]
    {\Large\color{sagegreen}\textbf{ANÁLISIS LÉXICO}}\\[0.3cm]
    {\large\color{papergray}Estudio de Palabras en Idiomas Originales}\\[1cm]
    
    \begin{tcolorbox}[colback=olivegreen!10, colframe=olivegreen, width=0.8\textwidth]
        \centering
        \textcolor{forestgreen}{\textbf{Proyecto: Análisis Exegético Profundo}}\\
        \textcolor{papergray}{Metodología de 8 Pasos para Estudio Bíblico Riguroso}
    \end{tcolorbox}
\end{center}

\vspace{1cm}

% Información del estudiante
\begin{infobox}{Información del Proyecto}
\begin{tabularx}{\textwidth}{|X|X|}
\hline
\textbf{Nombre del Estudiante:} & \underline{\hspace{6cm}} \\[0.8cm]
\hline
\textbf{Texto Bíblico:} & \underline{\hspace{6cm}} \\[0.8cm]
\hline
\textbf{Idioma Original:} & $\square$ Hebreo \quad $\square$ Griego \quad $\square$ Arameo \\[0.8cm]
\hline
\textbf{Fecha:} & \underline{\hspace{6cm}} \\[0.8cm]
\hline
\end{tabularx}
\end{infobox}

\section{Objetivo del Paso 2: Análisis Léxico}

El análisis léxico nos permite profundizar en el significado preciso de las palabras clave del texto en sus idiomas originales. Este paso es fundamental para entender los matices que pueden perderse en las traducciones.

\subsection{Descripción del Proceso}

Estudiaremos el significado de las palabras clave utilizando diccionarios griegos/hebreos, concordancias y herramientas lexicográficas. Examinaremos el rango semántico completo de cada término y cómo se usa en otros contextos bíblicos.

\begin{infobox}{Herramientas Recomendadas}
\textbf{Para Griego:} BDAG, Thayer, Vine's, Strong's Concordance\\
\textbf{Para Hebreo:} BDB, HALOT, TWOT, Strong's Concordance\\
\textbf{Software:} Logos, BibleWorks, Accordance, e-Sword (gratuito)
\end{infobox}

\section{Lista de Verificación}

\begin{checklistbox}
\textbf{Marca cada elemento completado:}

\begin{itemize}[leftmargin=1cm]
    \item[$\square$] Identificar 5-8 palabras clave del pasaje
    \item[$\square$] Buscar definiciones en diccionarios griegos/hebreos
    \item[$\square$] Examinar otros usos en el NT/AT usando concordancia
    \item[$\square$] Anotar matices y connotaciones específicas
    \item[$\square$] Considerar el campo semántico de cada palabra
    \item[$\square$] Evaluar diferentes traducciones posibles
    \item[$\square$] Documentar etimología cuando sea relevante
    \item[$\square$] Verificar el uso en literatura extrabíblica (LXX, etc.)
\end{itemize}
\end{checklistbox}

\newpage

\section{Análisis Detallado de Palabras Clave}

\subsection{Palabra Clave \#1}

\begin{lexiconbox}
\textbf{Palabra en Español:} \underline{\hspace{8cm}}\\[0.5cm]
\textbf{Palabra Original:} \underline{\hspace{4cm}} \quad \textbf{Transliteración:} \underline{\hspace{4cm}}\\[0.5cm]
\textbf{Strong's \#:} \underline{\hspace{3cm}} \quad \textbf{Versículo(s):} \underline{\hspace{5cm}}
\end{lexiconbox}

\textbf{Definición Principal (Diccionario):}
\vspace{2cm}

\textbf{Rango Semántico (Posibles significados):}
\begin{enumerate}[leftmargin=1cm]
    \item \underline{\hspace{12cm}}
    \item \underline{\hspace{12cm}}
    \item \underline{\hspace{12cm}}
    \item \underline{\hspace{12cm}}
\end{enumerate}

\textbf{Otros usos en la Escritura (mínimo 3 referencias):}

\vspace{0.5cm}

\begin{tabularx}{\textwidth}{|L{2.5cm}|L{5.5cm}|X|}
\hline
\textbf{Referencia} & \textbf{Contexto} & \textbf{Matiz/Significado} \\
\hline
\vspace{1.5cm} & \vspace{1.5cm} & \vspace{1.5cm} \\
\hline
\vspace{1.5cm} & \vspace{1.5cm} & \vspace{1.5cm} \\
\hline
\vspace{1.5cm} & \vspace{1.5cm} & \vspace{1.5cm} \\
\hline
\end{tabularx}

\vspace{0.5cm}

\textbf{Etimología y desarrollo histórico:}
\vspace{2cm}

\textbf{Significado más probable en este contexto:}
\vspace{2cm}

\newpage

\subsection{Palabra Clave \#2}

\begin{lexiconbox}
\textbf{Palabra en Español:} \underline{\hspace{8cm}}\\[0.5cm]
\textbf{Palabra Original:} \underline{\hspace{4cm}} \quad \textbf{Transliteración:} \underline{\hspace{4cm}}\\[0.5cm]
\textbf{Strong's \#:} \underline{\hspace{3cm}} \quad \textbf{Versículo(s):} \underline{\hspace{5cm}}
\end{lexiconbox}

\textbf{Definición Principal (Diccionario):}
\vspace{2cm}

\textbf{Rango Semántico (Posibles significados):}
\begin{enumerate}[leftmargin=1cm]
    \item \underline{\hspace{12cm}}
    \item \underline{\hspace{12cm}}
    \item \underline{\hspace{12cm}}
    \item \underline{\hspace{12cm}}
\end{enumerate}

\textbf{Otros usos en la Escritura (mínimo 3 referencias):}

\vspace{0.5cm}

\begin{tabularx}{\textwidth}{|L{2.5cm}|L{5.5cm}|X|}
\hline
\textbf{Referencia} & \textbf{Contexto} & \textbf{Matiz/Significado} \\
\hline
\vspace{1.5cm} & \vspace{1.5cm} & \vspace{1.5cm} \\
\hline
\vspace{1.5cm} & \vspace{1.5cm} & \vspace{1.5cm} \\
\hline
\vspace{1.5cm} & \vspace{1.5cm} & \vspace{1.5cm} \\
\hline
\end{tabularx}

\vspace{0.5cm}

\textbf{Etimología y desarrollo histórico:}
\vspace{2cm}

\textbf{Significado más probable en este contexto:}
\vspace{2cm}

\newpage

\subsection{Palabra Clave \#3}

\begin{lexiconbox}
\textbf{Palabra en Español:} \underline{\hspace{8cm}}\\[0.5cm]
\textbf{Palabra Original:} \underline{\hspace{4cm}} \quad \textbf{Transliteración:} \underline{\hspace{4cm}}\\[0.5cm]
\textbf{Strong's \#:} \underline{\hspace{3cm}} \quad \textbf{Versículo(s):} \underline{\hspace{5cm}}
\end{lexiconbox}

\textbf{Definición Principal (Diccionario):}
\vspace{2cm}

\textbf{Rango Semántico (Posibles significados):}
\begin{enumerate}[leftmargin=1cm]
    \item \underline{\hspace{12cm}}
    \item \underline{\hspace{12cm}}
    \item \underline{\hspace{12cm}}
    \item \underline{\hspace{12cm}}
\end{enumerate}

\textbf{Otros usos en la Escritura (mínimo 3 referencias):}

\vspace{0.5cm}

\begin{tabularx}{\textwidth}{|L{2.5cm}|L{5.5cm}|X|}
\hline
\textbf{Referencia} & \textbf{Contexto} & \textbf{Matiz/Significado} \\
\hline
\vspace{1.5cm} & \vspace{1.5cm} & \vspace{1.5cm} \\
\hline
\vspace{1.5cm} & \vspace{1.5cm} & \vspace{1.5cm} \\
\hline
\vspace{1.5cm} & \vspace{1.5cm} & \vspace{1.5cm} \\
\hline
\end{tabularx}

\vspace{0.5cm}

\textbf{Etimología y desarrollo histórico:}
\vspace{2cm}

\textbf{Significado más probable en este contexto:}
\vspace{2cm}

\newpage

\subsection{Palabras Adicionales (Continúa según necesidades)}

\begin{tabularx}{\textwidth}{|C{1.5cm}|L{3.5cm}|L{4cm}|X|}
\hline
\textbf{Palabra} & \textbf{Original/Strong's} & \textbf{Definición Clave} & \textbf{Significado Contextual} \\
\hline
\#4 & \vspace{1.5cm} & \vspace{1.5cm} & \vspace{1.5cm} \\
\hline
\#5 & \vspace{1.5cm} & \vspace{1.5cm} & \vspace{1.5cm} \\
\hline
\#6 & \vspace{1.5cm} & \vspace{1.5cm} & \vspace{1.5cm} \\
\hline
\#7 & \vspace{1.5cm} & \vspace{1.5cm} & \vspace{1.5cm} \\
\hline
\#8 & \vspace{1.5cm} & \vspace{1.5cm} & \vspace{1.5cm} \\
\hline
\end{tabularx}

\section{Síntesis Lexical}

\subsection{Campos Semánticos Identificados}

\textbf{¿Qué campos conceptuales emergen del vocabulario del texto?}
\vspace{2cm}

\textbf{¿Hay algún patrón en el uso del vocabulario por parte del autor?}
\vspace{2cm}

\subsection{Evaluación de Traducciones}

\begin{tabularx}{\textwidth}{|L{3cm}|L{3cm}|L{3.5cm}|X|}
\hline
\textbf{Palabra Original} & \textbf{Traducción Actual} & \textbf{Alternativas Posibles} & \textbf{Recomendación} \\
\hline
\vspace{1.5cm} & \vspace{1.5cm} & \vspace{1.5cm} & \vspace{1.5cm} \\
\hline
\vspace{1.5cm} & \vspace{1.5cm} & \vspace{1.5cm} & \vspace{1.5cm} \\
\hline
\vspace{1.5cm} & \vspace{1.5cm} & \vspace{1.5cm} & \vspace{1.5cm} \\
\hline
\vspace{1.5cm} & \vspace{1.5cm} & \vspace{1.5cm} & \vspace{1.5cm} \\
\hline
\end{tabularx}

\subsection{Implicaciones Teológicas}

\textbf{¿Qué nuevas perspectivas aporta este análisis léxico a tu comprensión del texto?}
\vspace{3cm}

\textbf{¿Hay algún matiz teológico importante que emerge del estudio de palabras?}
\vspace{3cm}

\begin{infobox}{Recordatorio Metodológico}
El significado de una palabra está determinado por su \textbf{contexto}, no solo por su definición de diccionario. Siempre considera el uso específico en el pasaje estudiado.
\end{infobox}

\section{Siguiente Paso}

Con el análisis léxico completado, procederás al \textbf{Paso 3: Análisis Gramatical}, donde examinaremos la estructura sintáctica y las relaciones gramaticales del texto.

\vspace{1cm}

\begin{center}
\begin{tcolorbox}[colback=mintgreen!20, colframe=mintgreen, width=0.7\textwidth]
\centering
\textcolor{forestgreen}{\textbf{¡Excelente Trabajo!}}\\
\textcolor{papergray}{Has profundizado en el significado original de las palabras clave.}
\end{tcolorbox}
\end{center}

\end{document}